% Reconstructed from the compiled lecture-note PDF.
\chapter{Curvature becomes a control variable}
\section{State-dependent curvatures}

\begin{displaytext}
Recall \
σj(t) = g(t)s∗j, g′ = gX, det X = 1.
\end{displaytext}

\begin{displaytext}
Differentiate twice: \
σ′j = gXs∗j, σ′′j = g(X2 + X′)s∗j.
\end{displaytext}

Since det g = 1, signed determinants are unchanged by g. For the three even-indexed curves define


\begin{displaytext}
κj = det(σ′2j, σ′′2j) = det Xs∗2j, (X2 + X′)s∗2j , j = 0, 1, 2.
\end{displaytext}

Using X−1 = −X, one can rewrite this as


\begin{displaytext}
κj = det s∗2j, (X + X−1X′)s∗2j .
\end{displaytext}

\begin{displaytext}
Convexity gives \
κj ≥0
\end{displaytext}

almost everywhere.


At least one of the three κj is positive at almost every time. Geometrically, if two vanish then two boundary curves are straight, and the third is a genuine hyperbolic arc with positive curvature. Algebraically, the star inequalities force the sum to be positive.


\section{Normalized curvature coordinates}

\begin{displaytext}
Set \
κ = κ0 + κ1 + κ2 > 0
\end{displaytext}

\begin{displaytext}
and define \
κj \
uj = . \
κ \
Then \
uj ≥0, u0 + u1 + u2 = 1.
\end{displaytext}

e2 = (0, 0, 1)


UT


13, 13, 13


\begin{displaytext}
e0 = (1, 0, 0) e1 = (0, 1, 0)
\end{displaytext}

each vertex means two curvatures vanish


\begin{figureplaceholder}
Figure 5.1: The triangular control set. Interior points distribute curvature among all three synchronized arcs; vertices concentrate it in one arc.
\end{figureplaceholder}

Thus the control lies in the standard two-simplex


\begin{displaytext}
n o UT = (u0, u1, u2) ∈R3 : uj ≥0, u0 + u1 + u2 = 1 .
\end{displaytext}

The normalization is crucial. The unnormalized curvatures depend on the speed of parametrization; the ratios uj do not.


\section{Encoding the control by a traceless matrix}

For each u ∈UT , define


\begin{displaytext}
( u1 −u2 u0 −2u1 −2u2 ) \
√ \
3 3 | Zu = | | u2 −u1 | ( u0 √ )∈sl2(R). \
3
\end{displaytext}

This affine map is chosen so that


\begin{displaytext}
uj = det(s∗2j, Zus∗2j), j = 0, 1, 2.
\end{displaytext}

Because three such determinant values determine a traceless 2 × 2 matrix, Zu is unique.


The vertices give three particularly simple matrices. We write


\begin{displaytext}
Z0 = Ze0, Z1 = Ze1, Z2 = Ze2.
\end{displaytext}

They are cyclically conjugate by R: RZjR−1 = Zj+1


with indices modulo 3.


\section{The star inequalities in control form}

A direct calculation gives


\begin{displaytext}
⟨Zu, X⟩= −2√ ρ2(X)u0 + ρ1(X)u1 + ρ0(X)u2 . \
3
\end{displaytext}

Since u ranges over the whole simplex,


\begin{displaytext}
ρ0, ho1, ho2 > 0 ⇐⇒ ⟨Zu, X⟩< 0 for every u ∈UT .
\end{displaytext}

The denominator appearing below is therefore never zero on an admissible trajectory.


\section{Deriving the differential equation for X}

From the definition of the curvature controls,


\begin{displaytext}
det s∗2j, (X + X−1X′)s∗2j = κuj = det(s∗2j, κZus∗2j).
\end{displaytext}

A traceless matrix is determined by these three determinant values, so


\begin{displaytext}
X + X−1X′ = κZu.
\end{displaytext}

\begin{displaytext}
Multiply by X: \
X′ = κXZu −X2.
\end{displaytext}

\begin{displaytext}
Since X2 = −I, this is \
X′ = κXZu + I. \
Take traces. Because tr X′ = 0, \
0 = κ tr(XZu) + 2,
\end{displaytext}

\begin{displaytext}
after using tr I = 2. Therefore \
2 \
κ = − ⟨X, Zu⟩.
\end{displaytext}

Substitute this back. The identity


XZu + ZuX = ⟨X, Zu⟩I


for traceless 2 × 2 matrices yields [Zu, X] X′ = ⟨Zu, X⟩.


This is the central state equation.


\begin{statementbox}{Theorem}
Theorem 5.1 (Lie-algebra state equation). The normalized body velocity satisfies
\end{statementbox}

\begin{displaytext}
[Zu, X] X′ = ⟨Zu, X⟩, u(t) ∈UT .
\end{displaytext}

Together with g′ = gX, this determines the boundary multicurve.


\section{Why the equation is isospectral}

Set Zu P = ⟨Zu, X⟩.


Then X′ = [P, X].


\begin{displaytext}
This is a Lax equation. Its basic property is that the characteristic polynomial of X is constant. For \
example, \
d \
tr(X2) = 2 tr(X[P, X]) = 2 tr(XPX −X2P) = 0. \
dt \
Thus det X is preserved, consistent with our normalization.
\end{displaytext}

If P were constant, the solution would be


X(t) = etP X(0)e−tP .


Even when P varies, the equation says that X(t) moves by infinitesimal conjugation on the determinant-one surface.


\section{What the vertices mean geometrically}

\begin{displaytext}
Suppose u = e0 = (1, 0, 0). Then
\end{displaytext}

\begin{displaytext}
κ1 = κ2 = 0, κ0 > 0.
\end{displaytext}

Two of the three even-indexed curves have zero curvature and hence move along straight lines. By the hyperbolic-envelope lemma, the remaining curve moves along a hyperbola with those lines as asymptotes. The same holds cyclically at the other vertices.


Therefore: vertex control ⇐⇒ two straight arcs and one hyperbolic arc .


A control that jumps among vertices produces a boundary that alternates among straight and hyperbolic pieces. This is why bang-bang control will imply a smoothed polygon.


\section{Interior and edge controls}

If u lies in the interior of UT , all three curvature components are positive. The circle corresponds to the center 1 1 1 u = 3, 3, 3 .


If u lies on an edge, exactly one curvature component vanishes. Such edge controls require special treatment because the Pontryagin maximization may fail to choose a unique endpoint of the edge. Chapter 10 develops an auxiliary edge-control problem to prove that true minimizers cannot slide arbitrarily along an edge.


\section{Reconstruction of the boundary}

\begin{displaytext}
Given measurable u(t) ∈UT and initial data g(0) = I, X(0) = X0 satisfying the star inequalities \
and det X0 = 1, solve \
[Zu, X] X′ = ⟨Zu, X⟩, g′ = gX.
\end{displaytext}

\begin{displaytext}
Then define \
σj(t) = g(t)s∗j.
\end{displaytext}

The state equation was derived precisely so that the normalized curvatures of these curves are uj. As long as the star inequalities and endpoint conditions hold, the multicurve reconstructs a balanced disk.


\begin{insightbox}{Geometric intuition}
\end{insightbox}

\begin{displaytext}
The control does not directly push the point g in the plane. It chooses how the total curvature \
is distributed among three synchronized boundary arcs. That curvature distribution changes \
the moving-frame velocity X; X then moves the frame g; and g generates all six boundary \
curves. \
u −→X′ −→g′ −→σj.
\end{displaytext}

\begin{insightbox}{Chapter summary}
\end{insightbox}

Convexity supplies three nonnegative curvature quantities. Dividing by their sum places the control in a triangle. The affine matrix Zu packages these curvature ratios. The complete normalized state dynamics is


\begin{displaytext}
g′ = gX, X′ = [Zu, X]/ ⟨Zu, X⟩.
\end{displaytext}

At a vertex of the control triangle, two boundary curves are straight and the third is hyperbolic. Consequently, finite vertex switching is exactly the geometric pattern of a smoothed polygon.

